\uuid{dkWg}
\titre{Hauteur du pic d'une densité normale}
\auteur{Maxime Nguyen}
\organisation{AMSCC}
\datecreate{2026-11-24}
\niveau{L2}
\module{Probabilité}
\chapitre{Probabilité continue}
\sousChapitre{Loi normale}
\difficulte{1}
\theme{loi normale, densité, écart-type}

\contenu{
\begin{question}
\qtype{onlyone}
On note $\mathcal{N}(\mu,\sigma)$ la loi normale d'espérance $\mu$ et d'écart-type $\sigma$ ; sa densité est $x \mapsto \frac{1}{\sigma\sqrt{2\pi}}\,e^{-\frac{(x-\mu)^2}{2\sigma^2}}$. Soit $X \sim \mathcal{N}(0;\,0{,}1)$. La valeur de la densité de $X$ en $x = 0$ est :
\begin{answers}
\bad{$\dfrac{1}{\sqrt{2\pi}}$}
\good{$\dfrac{10}{\sqrt{2\pi}}$}
\bad{$\dfrac{1}{10\sqrt{2\pi}}$}
\bad{un réel nécessairement inférieur ou égal à $1$, sans plus de précision}
\end{answers}
\begin{explanations}
En $x = 0$, la densité vaut
$\frac{1}{0{,}1\sqrt{2\pi}} = \frac{10}{\sqrt{2\pi}}$, qui est plus grand que
$1$ puisque $\sqrt{2\pi} < 10$ (il vaut environ $3{,}99$). A est la valeur pour
$\mathcal{N}(0,1)$ ($\sigma$ oublié) ; C multiplie par $\sigma$ au lieu de
diviser ; D repose sur l'idée fausse qu'une densité est majorée par 1.

Le maximum de la densité de $\mathcal{N}(\mu,\sigma)$ est
$\frac{1}{\sigma\sqrt{2\pi}}$ : l'aire totale reste égale à 1, donc plus
$\sigma$ est petit, plus la courbe est étroite et haute. Et
$\mathbb{P}(X = 0) = 0$.

Contre-sens visés : confondre la densité avec une probabilité, et en
déduire qu'elle est toujours inférieure à 1 (D) ; mal placer $\sigma$ dans la
formule de la densité (A, C).
\end{explanations}
\end{question}
}
